% Options for packages loaded elsewhere
% Options for packages loaded elsewhere
\PassOptionsToPackage{unicode}{hyperref}
\PassOptionsToPackage{hyphens}{url}
\PassOptionsToPackage{dvipsnames,svgnames,x11names}{xcolor}
%
\documentclass[
  12pt]{article}
\usepackage{xcolor}
\usepackage[margin=1in]{geometry}
\usepackage{amsmath,amssymb}
\setcounter{secnumdepth}{5}
\usepackage{iftex}
\ifPDFTeX
  \usepackage[T1]{fontenc}
  \usepackage[utf8]{inputenc}
  \usepackage{textcomp} % provide euro and other symbols
\else % if luatex or xetex
  \usepackage{unicode-math} % this also loads fontspec
  \defaultfontfeatures{Scale=MatchLowercase}
  \defaultfontfeatures[\rmfamily]{Ligatures=TeX,Scale=1}
\fi
\usepackage{lmodern}
\ifPDFTeX\else
  % xetex/luatex font selection
  \setmainfont[]{TeX Gyre Termes}
  \setmathfont[]{texgyretermes-math.otf}
\fi
% Use upquote if available, for straight quotes in verbatim environments
\IfFileExists{upquote.sty}{\usepackage{upquote}}{}
\IfFileExists{microtype.sty}{% use microtype if available
  \usepackage[]{microtype}
  \UseMicrotypeSet[protrusion]{basicmath} % disable protrusion for tt fonts
}{}
\usepackage{setspace}
\makeatletter
\@ifundefined{KOMAClassName}{% if non-KOMA class
  \IfFileExists{parskip.sty}{%
    \usepackage{parskip}
  }{% else
    \setlength{\parindent}{0pt}
    \setlength{\parskip}{6pt plus 2pt minus 1pt}}
}{% if KOMA class
  \KOMAoptions{parskip=half}}
\makeatother
% Make \paragraph and \subparagraph free-standing
\makeatletter
\ifx\paragraph\undefined\else
  \let\oldparagraph\paragraph
  \renewcommand{\paragraph}{
    \@ifstar
      \xxxParagraphStar
      \xxxParagraphNoStar
  }
  \newcommand{\xxxParagraphStar}[1]{\oldparagraph*{#1}\mbox{}}
  \newcommand{\xxxParagraphNoStar}[1]{\oldparagraph{#1}\mbox{}}
\fi
\ifx\subparagraph\undefined\else
  \let\oldsubparagraph\subparagraph
  \renewcommand{\subparagraph}{
    \@ifstar
      \xxxSubParagraphStar
      \xxxSubParagraphNoStar
  }
  \newcommand{\xxxSubParagraphStar}[1]{\oldsubparagraph*{#1}\mbox{}}
  \newcommand{\xxxSubParagraphNoStar}[1]{\oldsubparagraph{#1}\mbox{}}
\fi
\makeatother


\usepackage{longtable,booktabs,array}
\newcounter{none} % for unnumbered tables
\usepackage{calc} % for calculating minipage widths
% Correct order of tables after \paragraph or \subparagraph
\usepackage{etoolbox}
\makeatletter
\patchcmd\longtable{\par}{\if@noskipsec\mbox{}\fi\par}{}{}
\makeatother
% Allow footnotes in longtable head/foot
\IfFileExists{footnotehyper.sty}{\usepackage{footnotehyper}}{\usepackage{footnote}}
\makesavenoteenv{longtable}
\usepackage{graphicx}
\makeatletter
\newsavebox\pandoc@box
\newcommand*\pandocbounded[1]{% scales image to fit in text height/width
  \sbox\pandoc@box{#1}%
  \Gscale@div\@tempa{\textheight}{\dimexpr\ht\pandoc@box+\dp\pandoc@box\relax}%
  \Gscale@div\@tempb{\linewidth}{\wd\pandoc@box}%
  \ifdim\@tempb\p@<\@tempa\p@\let\@tempa\@tempb\fi% select the smaller of both
  \ifdim\@tempa\p@<\p@\scalebox{\@tempa}{\usebox\pandoc@box}%
  \else\usebox{\pandoc@box}%
  \fi%
}
% Set default figure placement to htbp
\def\fps@figure{htbp}
\makeatother


% definitions for citeproc citations
\NewDocumentCommand\citeproctext{}{}
\NewDocumentCommand\citeproc{mm}{%
  \begingroup\def\citeproctext{#2}\cite{#1}\endgroup}
\makeatletter
 % allow citations to break across lines
 \let\@cite@ofmt\@firstofone
 % avoid brackets around text for \cite:
 \def\@biblabel#1{}
 \def\@cite#1#2{{#1\if@tempswa , #2\fi}}
\makeatother
\newlength{\cslhangindent}
\setlength{\cslhangindent}{1.5em}
\newlength{\csllabelwidth}
\setlength{\csllabelwidth}{3em}
\newenvironment{CSLReferences}[2] % #1 hanging-indent, #2 entry-spacing
 {\begin{list}{}{%
  \setlength{\itemindent}{0pt}
  \setlength{\leftmargin}{0pt}
  \setlength{\parsep}{0pt}
  % turn on hanging indent if param 1 is 1
  \ifodd #1
   \setlength{\leftmargin}{\cslhangindent}
   \setlength{\itemindent}{-1\cslhangindent}
  \fi
  % set entry spacing
  \setlength{\itemsep}{#2\baselineskip}}}
 {\end{list}}
\usepackage{calc}
\newcommand{\CSLBlock}[1]{\hfill\break\parbox[t]{\linewidth}{\strut\ignorespaces#1\strut}}
\newcommand{\CSLLeftMargin}[1]{\parbox[t]{\csllabelwidth}{\strut#1\strut}}
\newcommand{\CSLRightInline}[1]{\parbox[t]{\linewidth - \csllabelwidth}{\strut#1\strut}}
\newcommand{\CSLIndent}[1]{\hspace{\cslhangindent}#1}



\setlength{\emergencystretch}{3em} % prevent overfull lines

\providecommand{\tightlist}{%
  \setlength{\itemsep}{0pt}\setlength{\parskip}{0pt}}



 


\makeatletter
\@ifpackageloaded{caption}{}{\usepackage{caption}}
\AtBeginDocument{%
\ifdefined\contentsname
  \renewcommand*\contentsname{Table of contents}
\else
  \newcommand\contentsname{Table of contents}
\fi
\ifdefined\listfigurename
  \renewcommand*\listfigurename{List of Figures}
\else
  \newcommand\listfigurename{List of Figures}
\fi
\ifdefined\listtablename
  \renewcommand*\listtablename{List of Tables}
\else
  \newcommand\listtablename{List of Tables}
\fi
\ifdefined\figurename
  \renewcommand*\figurename{Figure}
\else
  \newcommand\figurename{Figure}
\fi
\ifdefined\tablename
  \renewcommand*\tablename{Table}
\else
  \newcommand\tablename{Table}
\fi
}
\@ifpackageloaded{float}{}{\usepackage{float}}
\floatstyle{ruled}
\@ifundefined{c@chapter}{\newfloat{codelisting}{h}{lop}}{\newfloat{codelisting}{h}{lop}[chapter]}
\floatname{codelisting}{Listing}
\newcommand*\listoflistings{\listof{codelisting}{List of Listings}}
\makeatother
\makeatletter
\makeatother
\makeatletter
\@ifpackageloaded{caption}{}{\usepackage{caption}}
\@ifpackageloaded{subcaption}{}{\usepackage{subcaption}}
\makeatother
\usepackage{bookmark}
\IfFileExists{xurl.sty}{\usepackage{xurl}}{} % add URL line breaks if available
\urlstyle{same}
\hypersetup{
  pdftitle={Tax Evasion and Productivity: Identifying and Correcting for Cost Overreporting},
  pdfauthor={Hans Martinez},
  pdfkeywords={Tax Evasion, Cost Overreporting, Production Function
Estimation, Productivity},
  colorlinks=true,
  linkcolor={black},
  filecolor={Maroon},
  citecolor={black},
  urlcolor={blue},
  pdfcreator={LaTeX via pandoc}}


\title{Tax Evasion and Productivity: Identifying and Correcting for Cost
Overreporting\thanks{I thank Salvador Navarro, David Rivers, and Timothy
Conley for helpful comments. All errors are my own.}}
\author{Hans Martinez\thanks{Department of Economics, Western
University. Email: \href{mailto:hmarti33@uwo.ca}{hmarti33@uwo.ca}. Website: \href{https://hansmartinez.com}{hansmartinez.com}.}}
\date{October 08,
2026\\[0.5em]\href{https://hansmartinez.com/project/tax\_prod/}{\small\textbf{Latest version, click here}}}
\begin{document}
\maketitle
\begin{abstract}
Value-added taxes (VAT) allow firms to deduct the taxes paid on
purchases from the taxes collected on sales, creating an incentive for
firms to overreport the purchases of inputs to reduce their tax bill.
The fundamental issue in detecting input overreporting stems from two
unobservables, true inputs and productivity. Because firms differ in
productivity, a firm reporting high input use, for a given level of
output, may be either overreporting or simply less productive. To
distinguish between overreporting and productivity, I use a structural
production-function approach and a benchmark group, a subset of firms
that report inputs correctly. Because the benchmark firms are assumed to
share the same technology as the other firms, I estimate the
technology's parameters from them. Using these estimates, I compare what
potentially tax-evading firms report with predictions of the true inputs
they would use. I apply the method to a Colombian manufacturing survey
from 1981 to 1991. During this period, the Colombian sales tax worked as
a VAT and the data include a subset of firms that, due to government and
market scrutiny, reported inputs correctly.

My approach yields a simple and robust statistical test to detect input
overreporting. Under the null hypothesis of correct reporting, the test
is agnostic about the specification of both the probability of detection
and the cost of overreporting. Among the 5 largest industries, the test
rejects the null hypothesis in all 4 industries subject to the sales
tax, and, as expected, it fails to reject in the industry exempt from
the sales tax. In these 4 industries, firms overreport on average 14.7
percent of their true inputs in the median industry. Exploiting a 1983
fiscal reform in Colombia that raised the sales tax rate, I then
document that firms overreport more after the increase. The response
implies a midpoint elasticity of overreporting with respect to the tax
rate of 5.9.

The magnitude of detected overreporting matters for methods of
estimating production functions that exploit a flexible input to
identify productivity (Levinsohn and Petrin 2003; Doraszelski and
Jaumandreu 2013; Ackerberg et al. 2015; Gandhi et al. 2020). These
methods typically assume that materials are flexible, and their demand
or first-order conditions reveal productivity. However, materials are
deductible under the VAT. Therefore, firms have the incentive to
overreport the input these methods rely on. I show that failing to
correct for overreporting yields larger output elasticities of
materials, and more dispersed and less persistent productivity
distributions. Using the production-function parameters, I estimate a
structural model of input overreporting. With the model, I study how
government revenue changes as expected deductions, true and fictitious,
respond to the tax rate. At the baseline rate, a 1 percent increase in
the sales tax rate on purchases raises expected deductions by about 1.4
percent: firms overreporting more add 0.2 points, and firms buying more
materials add another 0.2. At the current rate, the government pays
potentially tax-evading firms net refunds, because their claims exceed
the sales tax they owe on their sales. Each 1 percent increase in the
rate lowers net revenue by about 1.7 percent of that sales tax, and a
cut of about 17 percent would turn the refunds into positive revenue.
Undetected overreporting costs the government 4.7 percent of the sales
tax firms in tax evading industries owe on their sales.
\end{abstract}

\begin{center}
\textbf{JEL Classification:} H26, D22, L60, O47 \\
\textbf{Keywords:} Tax Evasion, Cost Overreporting, Production Function Estimation, Productivity
\end{center}
\bigskip


\setstretch{1.15}
\section*{References}\label{references}
\addcontentsline{toc}{section}{References}

\protect\phantomsection\label{refs}
\begin{CSLReferences}{1}{1}
\bibitem[\citeproctext]{ref-Ackerberg2015}
Ackerberg, Daniel A, Kevin Caves, and Garth Frazer. 2015.
{``Identification Properties of Recent Production Function
Estimators.''} \emph{Econometrica} 83 (6): 2411--51.
\url{https://doi.org/10.3982/ECTA13408}.

\bibitem[\citeproctext]{ref-Doraszelski2013}
Doraszelski, Ulrich, and Jordi Jaumandreu. 2013. {``{R\&D} and
Productivity: Estimating Endogenous Productivity.''} \emph{The Review of
Economic Studies} 80 (4): 1338--83.
\url{https://doi.org/10.1093/restud/rdt011}.

\bibitem[\citeproctext]{ref-Gandhi2020}
Gandhi, Amit, Salvador Navarro, and David A Rivers. 2020. {``On the
Identification of Gross Output Production Functions.''} \emph{Journal of
Political Economy} 128 (8): 2973--3016.
\url{https://doi.org/10.1086/707736}.

\bibitem[\citeproctext]{ref-Levinsohn2003}
Levinsohn, James, and Amil Petrin. 2003. {``Estimating Production
Functions Using Inputs to Control for Unobservables.''} \emph{Review of
Economic Studies} 70 (2): 317--41.
\url{https://doi.org/10.1111/1467-937X.00246}.

\end{CSLReferences}




\end{document}
